% Topic proposal: original template format, five mandatory sections.
\documentclass[11pt]{article}
\usepackage[utf8]{inputenc}
\usepackage[T1]{fontenc}
\usepackage[margin=1in]{geometry}
\usepackage{amsmath,amssymb}
\usepackage{xcolor}
\usepackage[round]{natbib}
\usepackage[colorlinks=true,allcolors=blue!50!black]{hyperref}
\title{Public Dollar Backing and Voluntary Bank Buffers}
\author{Carlos G\'omez Puic\'an\\
\small Universidad del Pac\'ifico\\
\small\url{https://github.com/carlosgomezp28/ai-project}}
\date{AI \& Economic Modeling: Track B, October 2026}
\begin{document}
\maketitle

\section{The question and Track B}
How does additional public dollar backing change a representative bank's optimal voluntary dollar buffer, conditional on aggregate liquidity conditions? This is Track B: a private decision margin within a thesis framework. For motivation, \citet[Box~7]{bcrp2024} attributes substantially higher dollar reserve requirements than soles requirements to higher foreign-currency liquidity risk. This motivates distinguishing currencies; it does not establish our comparative static.

The baseline is joint IE1 work by Carlos G\'omez Puic\'an and Fabi\'an M\'endez Su\'arez, supervised by Marco Ortiz, used here with Fabi\'an's permission. It already anticipates endogenous voluntary buffers and a future fixed point. Carlos's individual contribution completes the private optimization closure, introduces reduced-form private costs, derives its FOC, and characterizes a conditional best response to dollar backing.

\section{The baseline and the model the thesis lacked}
The inherited static economy has two representative classes $k\in\{L,S\}$, with masses $\mu_k$, differing in idiosyncratic liquidity-shock dispersion and interbank access, and currencies $c\in\{s,u\}$ (soles, dollars). Deposits satisfy $d_k^u=w_kd_k$ and $d_k^s=(1-w_k)d_k$, with $w_k$ fixed. Equity allocation satisfies $\sum_cE_k^c=E_k$. Define $v_k^c:=u_k^c$ in the joint IE1 notation to separate the buffer symbol from the dollar superscript. Required reserves $\rho^c$, voluntary buffers $v_k^c$, and total reserves $m_k^c=\rho^c+v_k^c$ are ratios to deposits; credit $b_k^c$ and deposits are amounts.

The bank chooses its portfolio before the shock, then accesses interbank funding and public coverage. The inherited shock is $\omega_k^c=\beta_kA^c+e_k^c$, with $e_k^c\sim U[-a_k^c,a_k^c]$, $A^s=0$, and $A^u\in\{-\eta,0\}$ with probabilities $p,1-p$. We maintain this conditional uniform law, $a_k^c>0$, $\eta>0$, and $0\le p\le1$. In NU, required reserves are not usable and back a collective fund; in U they are usable and their principal is not counted twice. For $r\in\{NU,U\}$, with $\chi_{NU}=0$, $\chi_U=1$, normalized liquidity is
\[
 x_k^c:=s_k^c/d_k^c=\beta_kA^c+e_k^c+v_k^c+\chi_r\rho^c.
\]
Inherited expected deficit and surplus amounts are $S_k^{c-}(A)=\mathbb E_e[(-s_k^c)^+\mid A]$ and $S_k^{c+}(A)=\mathbb E_e[(s_k^c)^+\mid A]$; $\pi_k^{c-}(A)=\Pr(s_k^c<0\mid A)$.

\newpage
Interbank contact has probability $\tau_k=1-e^{-\lambda_k}$, $\lambda_k\ge0$. Admitted supply and demand and residual deficits are
\begin{align*}
 P^c(A)&=\sum_k\mu_k\tau_kS_k^{c+}(A),&
 Q^c(A)&=\sum_k\mu_k\tau_kS_k^{c-}(A),\\
 \Psi_k^{-,c}(A)&=\tau_k\min\{1,P^c(A)/Q^c(A)\},&
 \Delta^c(A)&=\sum_k\mu_k(1-\Psi_k^{-,c}(A))S_k^{c-}(A).
\end{align*}
Dollar capacity and public coverage are
\[
 F^u=\mathbf1_{\{NU\}}\nu_F^u\rho^uD^u+R_D,
 \qquad \phi^u(A)=\min\{1,F^u/\Delta^u(A)\},
 \qquad D^u=\sum_k\mu_kd_k^u,
\]
where $\nu_F^u\in[0,1]$ is the mobilizable fraction and $R_D\ge0$ additional dollar backing. We write $\phi^u$ for inherited $\phi$, and take $\phi^s=1$ under the elastic soles benchmark. Displayed ratios require positive denominators; no-deficit conventions remain to be specified for numerical cases.

The joint proposal envisages optimizing buffers but lacks the following explicit private-cost objective. Conditional on aggregate coverage schedules, the representative bank minimizes cost per unit of deposits:
\begin{equation}\label{eq:problem}
 \min_{0\le v_k^c\le\bar v_k^c}V_k^c(v_k^c)
 =\kappa_k^c v_k^c+\mathbb E_{A^c}[h_k^c(A^c)L_k^c(v_k^c,A^c)],
\end{equation}
\begin{equation}\label{eq:h}
 h_k^c(A)=\Psi_k^{-,c}(A)q_I^c+
 (1-\Psi_k^{-,c}(A))\{\phi^c(A)q_F^c+(1-\phi^c(A))\ell^c\}.
\end{equation}
The new primitives are marginal buffer cost $\kappa_k^c>0$ and nonnegative unit costs $q_I^c,q_F^c,\ell^c$ for interbank funding, public funding, and uncovered deficits. The exclusive weights in \eqref{eq:h} sum to one. These costs are reduced-form, not identified with contractual interest rates. Imposing the \emph{new} restriction $b_k^c\ge0$ on inherited $b_k^c+d_k^c(\rho^c+v_k^c)=d_k^c+E_k^c$ yields $\bar v_k^c=1+E_k^c/d_k^c-\rho^c$, maintained positive with $d_k^c>0$ and fixed equity allocation.

\section{The result expected}
Define the new normalized object $L_k^c(v,A):=S_k^{c-}(v,A)/d_k^c$. Suppressing indices, put $z=\beta A+v+\chi_r\rho$. Under the conditional uniform law,
\[
 L(v,A)=\begin{cases}-z,&z\le-a,\\(a-z)^2/(4a),&-a<z<a,\\0,&z\ge a.\end{cases}
 \qquad L_v=-\pi^-(v,A).
\]
An extra buffer unit reduces shortfall by one only in deficit realizations. The pieces join in value and first derivative. For the amount, $S_v^-=-d\pi^-$; the deposit factor cannot be omitted. With $h$ independent of the individual choice in this conditional exercise, differentiation of \eqref{eq:problem} gives the interior FOC
\begin{equation}\label{eq:foc}
 \kappa_k^c=\mathbb E_{A^c}[h_k^c(A^c)\pi_k^{c-}(v_k^{c,*},A^c)].
\end{equation}
The objective is convex. A regular interior solution requires $0<v^*<\bar v$, local smoothness away from support and coverage kinks, and curvature $\mathcal H:=V_{vv}(v^*,R_D)>0$. Positive curvature holds if a positive-weight state has $h>0$ and its threshold strictly inside the uniform support. Positive deficit probability alone is insufficient. At the boundaries the convex problem instead uses $V_v(0)\ge0$ or $V_v(\bar v)\le0$.

For dollars, freeze $P^u(A),Q^u(A),\Delta^u(A)$ across backing scenarios, not across states; also freeze the other primitives, shock law, regime, and feasible set. In a rationed state, $\Delta^u(A)>0$ and $F^u<\Delta^u(A)$ imply
\[
 h_{R_D}=-\frac{(1-\Psi_k^{-,u})(\ell^u-q_F^u)}{\Delta^u},
 \qquad \frac{dv_k^{u,*}}{dR_D}
 =\frac{\mathbb E_A[\pi_k^{u-}h_{R_D}]}{\mathcal H}.
\]
Thus the candidate strict sign is $dv_k^{u,*}/dR_D<0$ at a regular interior optimum if $\ell^u>q_F^u$ and at least one positive-probability state simultaneously has a deficit probability greater than zero, incomplete interbank coverage $\Psi_k^{-,u}<1$, and rationed public coverage responsive to $R_D$. Saturated states contribute zero locally. This is a conditional individual best response, \emph{not an equilibrium theorem}; aggregate feedback, corners, and coverage kinks lie outside this derivative claim.

\section{Why this is not already solved}
\textbf{A. Internal gap.} The joint IE1 framework anticipates endogenous self-insurance and an equilibrium fixed point, but does not specify \eqref{eq:problem}--\eqref{eq:h} or derive \eqref{eq:foc}. The contribution is this private closure and its conditional characterization, not the introduction of endogenous buffers or sole authorship of the baseline.

\textbf{B. External positioning.} The general mechanism already exists. \citet{ratnovski2009} studies precautionary bank liquidity choices and distortions from expected lender-of-last-resort support; our reduced-form closure does not claim to discover that incentive. \citet{farhi2012} link expected public intervention to private balance-sheet choices and collective moral hazard; we hold aggregate conditions fixed and do not establish their strategic or systemic result. \citet{freixas2011} study liquidity shocks, interbank allocation, and central-bank policy; we characterize the private buffer margin within the specified IE1 coverage rules. \citet{rochet2004} analyze failures of private liquidity provision and lender-of-last-resort intervention; we do not claim a coordination-failure theorem.

The defensible contribution is a closure/application within the joint framework combining two bank classes, soles and dollars, voluntary and required reserve ratios, interbank coverage, U/NU usability, and limited dollar backing. Initial source checks support these broad comparisons. Exact proposition comparisons, appendices, working-paper versions, and citing papers remain to be reviewed; literature-wide novelty is uncertain and is not claimed.

\section{Plan and risks}
Prove the piecewise deficit derivative, characterize corners and regular interior solutions, and verify the conditional backing sign. The existing script has 28 manually passing symbolic checks; it verifies conditional algebra, not equilibrium existence, calibration, or contractual accounting. Numerically illustrate interior, zero-buffer, and coverage-saturation cases. Then study the aggregate fixed point by recomputing matching and coverage consistently with buffers; report existence or convergence limitations rather than treating numerical convergence as a proof. Formalize the numbered final-paper results in Lean, disclose added hypotheses or open gaps, and keep the paper, handwritten derivations, and checks aligned.

The main risk is that aggregate feedback overturns the conditional sign, or the optimum reaches a corner or coverage kink. The fallback is an honestly local/conditional result with explicit boundaries. Literature positioning, fixed-point analysis, and parameter interpretation remain work in progress; no net-transfer or welfare claim follows from this private problem.

\bibliographystyle{plainnat}
\bibliography{../paper/references}
\end{document}
